\documentclass{jsarticle}
\usepackage[pdftex]{graphicx}
\usepackage{mathtools,amsmath,amssymb}
\usepackage{multirow}
\usepackage{xcolor}
\usepackage{bm}
\usepackage[dvipdfmx]{hyperref}

\def\be#1\ee{\begin{align}#1\end{align}}


\begin{document}

-----------------------
山本さん（一部）
-----------------------\\

---数値計算について---\\
\textcolor{blue}{了解です、行列が結構小さくて完全に対角化できるのですね。500umだと思っていたので、セルサイズを100nmにとってもN=5000で、固有値を全て計算およびプロットするのは時間がかかるかと考えての質問でした。ちなみにトータル6Nある固有値の内半分は正、残りは負で、反転対称性があるので基本的に正負のペアは同じ周波数を持っていると思いますが、正のモードだけプロットしているという理解でよろしいでしょうか?}\\

紛らわしいミスをしてしまい申し訳ありませんでした。。。

プロットについて、ご推察の通りで、正負の固有値が対称に出現し、そのうち正のものだけをプロットしています。\\


------
\textcolor{blue}{そうなのでしょうか？式(8)を見ると$M_{0x}=0$であれば任意の変形に対して$G_u=0$となって音波のモードによらず結合は生じないように見えます。対称性の一般論で結合があって良くてもモデルの特殊性によって実際には結合が出ないケースは多く、慎重な考察が必要だと思います。}\\

ご指摘ありがとうございます。結論から申し上げますと、ご指摘の通り、今回のモデル・理論の範疇ではその特殊性によりcouplingは起こらないが、対称性によって結合が禁止されているわけではない。ということだと思います。

まず、対称性の議論としては、前回のreplyでお答えした通り、$M_y$固有値が$-1$の$u_y$のモードはMSWと結合できるはずです。つまり、一般にはshear horizontal waveとMSWはcoupleするはずです。

一方、今回の理論では磁場を決める式(3),(4)には、$y$方向について一様であるという仮定のもとで$u_y$が現れないことから、外部磁場の向きがどうであれ、shear horizontal waveとMSWは結合しないです。この点が系の(今回考えている静磁近似$+$線形という理論の)特殊性です。

原稿で議論しているLamb波との結合については、Lamb波である時点で$M_y$固有値は$+1$であり($(u_x,u_z)\to(u_x,u_z)$)、MSWは$M_{0,x}=0$ならば必ず$M_y$固有値が$-1$なので($(m_x,m_z)\to(-m_x,-m_z)$)、この対称性セクターの違いにより$G^u=0$になっていると理解できると思います。\\

まとめると、
\begin{itemize}
    \item Lamb波$u_x,u_z$との不結合$\to\ \because$鏡映固有値が$\varphi=\pi/2$では異なる
    \item Shear horizontal波$u_y$との不結合$\to\ \because$静磁近似と線形近似により任意の$\varphi$で$u_y$が効かない
\end{itemize}
ということだと思います。原稿では$u_y$に触れていなかったので(11)式の下の時点で$u_y$とは結合しないことに触れ、Fig.2 (c)の説明で、「Lambは$M_y$により結合しない。$u_y$は対称性的にはその限りではないが、今回の理論の範疇ではそもそも結合しない。」といったことを追加しました。\\


---磁気弾性結合定数について---\\
丁寧にご教示いただきありがとうございます。大変勉強になりました。完全に受け売りになりますが$g_{\mathrm{eff}}$の議論に続けて磁気弾性結合定数についての議論も追記いたしました。


-----------------------
藤本さん
-----------------------\\

---イントロ二段落目について---\\
\textcolor{blue}{''In this Letter, we extend the ..."という書き方だと、Kalinikos-Slavinの単なる拡張としてとらえられてしまう可能性がある。}\\

文頭でまず何をしたかの主張をするように書き換え、本論文の主要結果である「計算方法の提供」に焦点を合わせました。Kalinikos--Slavinのextensionというのをやめて、``along the same line as..."という表現にしてみました。\\


その他、軽微な修正をしました。ご指摘ありがとうございました。

% \begin{thebibliography}{99}
% \bibitem{lambWavesElasticPlate1917}
% H.~Lamb, On Waves in an Elastic Plate, \href{https://doi.org/10.1098/rspa.1917.0008}{Proc. A \textbf{93}, 114 (1917)}.

% \end{thebibliography}

\end{document}